\documentclass[11pt,a4paper]{article}
\usepackage[T1]{fontenc}
\usepackage[utf8]{inputenc}
\usepackage{amsmath,amssymb,xcolor}
\usepackage[margin=25mm]{geometry}
\usepackage[colorlinks=true,urlcolor=blue]{hyperref}
\definecolor{nexorablue}{HTML}{2459B5}
\setlength{\parindent}{0pt}
\setlength{\parskip}{7pt}
\setlength{\emergencystretch}{2em}
\renewcommand{\familydefault}{\sfdefault}
\begin{document}

{\small\bfseries\color{nexorablue} NEXORA CPGE}\par
{\small MP / NEX-STR-02 / Corrigé}\par
{\LARGE\bfseries Ordre et image d’un sous-groupe cyclique}\par
\vspace{0.25em}\hrule\vspace{0.6em}

\section*{Ordre}
On a $30\mid12n$ si et seulement si $5\mid2n$, donc exactement si $5\mid n$. L’ordre est $5$ et le noyau du morphisme est $5\mathbb Z$.

\section*{Image}
Les cinq classes obtenues sont $\overline0$, $\overline{12}$, $\overline{24}$, $\overline6$ et $\overline{18}$. Ce sont précisément les classes multiples de $\overline6$.


\end{document}
